% Generado por scripts/70_tablas_figuras.R a partir de resultados-v1. No editar a mano.
\begin{table}[htbp]
\centering
\small
\caption{Tabla 6.1—Efecto de los proyectos REDD+ sobre la deforestación municipal}
\label{tab:tabla_6_1}
\begin{threeparttable}
\begin{tabular}{@{}lcccc@{}}
\toprule
 & (1) & (2) & (3) & (4) \\
 & Principal & Sin SINAP & Sin covariables & D1 original \\
\midrule
ATT agregado & $-$152,5 & $-$169,4 & $-$183,3 & $-$163,7 \\
 & (186,5) & (194,5) & (298,3) & (168,3) \\
 & [$-$518; 213] & [$-$550; 212] & [$-$768; 401] & [$-$493; 166] \\
\midrule
ATT balanceado, e = 0--5 & $-$200,1 & — & $-$247,3 & $-$202,9 \\
 & (196,9) &  & (304,9) & (189,2) \\
Efecto relativo a la media previa & $-$15,1 \% & $-$16,8 \% & $-$18,2 \% & $-$15,9 \% \\
Media previa de los tratados (ha/año) & 1.009,5 & 1.009,5 & 1.009,5 & 1.027,7 \\
\midrule
Covariables & 11 & 10 & No & 11 \\
Muestra de tratados & D1 depurada & D1 depurada & D1 depurada & D1 original \\
Municipios tratados & 36 & 36 & 36 & 41 \\
Municipios de control & 986 & 986 & 986 & 987 \\
\bottomrule
\end{tabular}
\begin{tablenotes}[flushleft]
\footnotesize
\item \textit{Notas:} Estimador de Callaway y Sant’Anna (2021) doblemente robusto, con grupo de comparación de municipios aún no tratados y cohortes de adopción 2013--2019. Variable dependiente: pérdida anual de cobertura arbórea (ha). Errores estándar entre paréntesis, por bootstrap multiplicador agrupado por municipio (999 repeticiones); intervalos de confianza al 95 \% entre corchetes. La columna (1) usa las 11 covariables de línea base (D4); la (2) excluye el área en el SINAP. La D1 depurada excluye seis municipios cuyo único proyecto REDD+ fue retirado o rechazado. Controles limpios: municipios sin ningún vínculo con proyectos de carbono. El efecto relativo divide el ATT agregado por la media previa de los tratados.
\item \textit{Fuente:} Cálculos propios con Global Forest Change v1.12 y los registros Verra, RENARE y Cercarbono.
\end{tablenotes}
\end{threeparttable}
\end{table}
